% ==========================================================================
\section{Language A: commutative diagrams}
\label{sec:catdiag}

\paragraph{Marks.} Dots (objects) and arrows (morphisms); an arrow may be labelled.
\paragraph{Legal moves.} \emph{Pasting}: two commuting cells sharing an edge form a
commuting region. \emph{Filling}: if all inner cells commute, the outer boundary commutes.
\paragraph{Soundness.} Associativity and unitality of composition; this is a theorem, and
it is why the language is safe.

\subsection{The inverse, in four arrows}

``$g$ is inverse to $f$'' is by definition a statement about two triangles.

\begin{figure}[H]
\centering
\begin{tikzcd}[column sep=large,row sep=large]
X \arrow[r,"f"] \arrow[dr,"\id_X"'] & Y \arrow[d,"g"] \\
 & X
\end{tikzcd}
\qquad
\begin{tikzcd}[column sep=large,row sep=large]
Y \arrow[r,"g"] \arrow[dr,"\id_Y"'] & X \arrow[d,"f"] \\
 & Y
\end{tikzcd}
\qquad
\begin{tikzcd}[column sep=large,row sep=large]
Y \arrow[r,"g'"] \arrow[dr,"\id_Y"'] & X \arrow[d,"f"] \arrow[r,"g"] & Y \\
 & Y \arrow[ur,"g"'] &
\end{tikzcd}
\caption{Left and centre: the two triangles that say $g=f^{-1}$. Right: pasting the
left-inverse triangle of $g'$ onto the arrow $g$ shows $g'=g'\circ\id=g'\circ(f\circ g)
=(g'\circ f)\circ g=\id\circ g=g$. The inverse is unique -- drawn, not computed.
Certified in \lean{DiagramProofs.inverse\_unique} and, in an arbitrary category, in
\lean{SyllabusOverlap.inverse\_unique}.}
\label{fig:invunique}
\end{figure}

One picture covers both halves of the syllabus overlap:
\begin{itemize}\itemsep2pt
\item \emph{Kalkulus}, exercises 2.3--2.9 of \lean{kalkulus\_gyakorlo.pdf}: an answer
  $f^{-1}$ is right precisely when \emph{both} triangles close. Students routinely check
  only one; the diagram makes the omission visible.
\item \emph{Line\'aris algebra}: $MN=I$ with $\det M\neq 0$ forces $N=M^{-1}$ -- the same
  figure with $X=Y=\R^n$ and arrows the linear maps.
\end{itemize}

\subsection{The triangle inequality as composition (Lawvere)}

The triangle inequality looks like a statement about numbers, not a diagram. Lawvere's
observation is that it \emph{is} a composition law. Read a metric space as a category
whose objects are the points, whose hom-object from $x$ to $y$ is the number $\dd(x,y)$,
whose composition is $+$, and in which ``$\ge$'' is the direction of implication.

\begin{figure}[H]
\centering
\begin{tikzcd}[column sep=huge,row sep=large]
x \arrow[r,"\dd(x{,}y)"] \arrow[dr,"\dd(x{,}z)"'] & y \arrow[d,"\dd(y{,}z)"] \\
 & z
\end{tikzcd}
\qquad\raisebox{10mm}{$\rightsquigarrow$}\qquad
\begin{tikzpicture}[baseline=(current bounding box.center)]
  \coordinate (x) at (0,0);
  \coordinate (y) at (2.4,1.3);
  \coordinate (z) at (3.5,-0.5);
  \draw[vec] (x) -- node[above left]{$\dd(x,y)$} (y);
  \draw[vec] (y) -- node[right]{$\dd(y,z)$} (z);
  \draw[vec,dashed] (x) -- node[below]{$\dd(x,z)$} (z);
  \fill (x) circle (1.6pt) node[left]{$x$};
  \fill (y) circle (1.6pt) node[above]{$y$};
  \fill (z) circle (1.6pt) node[right]{$z$};
\end{tikzpicture}
\caption{The same statement twice. Left: the composite arrow $x\to z$ exists \emph{at
cost} $\dd(x,y)+\dd(y,z)$, while the hom-object $\dd(x,z)$ records the cheapest cost;
hence $\dd(x,z)\le\dd(x,y)+\dd(y,z)$. Right: the schoolbook triangle. Certified in
\lean{DiagramProofs.detour}.}
\label{fig:lawvere}
\end{figure}

Two things become visible in the categorical drawing that the schoolbook triangle hides.

\begin{figure}[H]
\centering
\begin{tikzpicture}[x=1cm,y=1cm]
  \foreach \i/\lab in {0/{f_0},1/{f_1},2/{f_2},3/{f_3},4/{f_4}}
    {\fill (1.5*\i,0) circle (1.6pt) node[below=3pt]{$\lab$};}
  \foreach \i in {0,1,2,3}
    {\draw[vec] (1.5*\i+0.14,0) -- (1.5*\i+1.36,0);}
  \draw[vec,dashed] (0,0.8) .. controls (3,1.9) .. (6,0.8);
  \node at (3,2.0) {$\dd(f_0,f_4)$};
  \node at (0.75,0.22) {$\scriptstyle \dd_0$};
  \node at (2.25,0.22) {$\scriptstyle \dd_1$};
  \node at (3.75,0.22) {$\scriptstyle \dd_2$};
  \node at (5.25,0.22) {$\scriptstyle \dd_3$};
\end{tikzpicture}
\caption{Composing a \emph{chain}: iterate the pasting move $n-1$ times to get
$\dd(f_0,f_n)\le\sum_{i<n}\dd(f_i,f_{i+1})$. With $f_k=\sum_{i<k}a_i$ in $\R$ this is
exercise 1.13 (generalised triangle inequality); with $f_k=\sum_{i<k}v_i$ in $\R^n$ it is
$\norm{\sum v_i}\le\sum\norm{v_i}$. One drawing, two courses. Certified in
\lean{TriangleOverlap.dist\_chain} and \lean{DiagramProofs.abs\_sum\_range\_le\_sum\_abs}.}
\label{fig:chain}
\end{figure}

\begin{figure}[H]
\centering
\begin{tikzpicture}[scale=1.0]
  \draw[thick] (0,0) circle (2.1);
  \node at (-1.55,1.55) {$B(x,r)$};
  \fill (0,0) circle (1.6pt) node[below left]{$x$};
  \coordinate (y) at (1.1,0.35);
  \draw[thick,dashed] (y) circle (0.95);
  \fill (y) circle (1.6pt) node[above]{$y$};
  \draw[vec,gray!70!black] (0,0) -- node[below,pos=0.55,font=\scriptsize]{$\dd(x,y)$} (y);
  \draw[vec,gray!70!black] (y) -- ++(0.67,0.67)
        node[right,font=\scriptsize]{$r-\dd(x,y)$};
\end{tikzpicture}
\caption{The triangle inequality drawn as an \emph{inclusion of discs}: if
$\dd(x,y)<r$ then $B\bigl(y,\,r-\dd(x,y)\bigr)\subseteq B(x,r)$. This is the picture
every ``open set'' and every limit argument in the course actually uses; drawing the two
nested circles is equivalent to writing the inequality. Certified in
\lean{DiagramProofs.ball\_subset\_ball\_of\_mem}.}
\label{fig:balls}
\end{figure}

\begin{recipe}[Draw an estimate as an arrow]
Whenever the quantity you must bound is the ``cost of getting from $A$ to $B$'', draw $A$
and $B$ as dots and every estimate you already have as a labelled arrow. Chaining arrows
adds costs; the bound you want is the cheapest path. Every $\varepsilon/3$ argument in
\emph{Kalkulus I} is this picture with three arrows:
\[
|f(x)-L|\;\le\;\underbrace{|f(x)-f_n(x)|}_{\varepsilon/3}
+\underbrace{|f_n(x)-L_n|}_{\varepsilon/3}+\underbrace{|L_n-L|}_{\varepsilon/3}.
\]
\end{recipe}

\subsection{The bridge between the two courses, in one arrow}

The two syllabus items are joined by a single morphism of the category $\Met$: the length
map $\norm{\cdot}\colon\R^n\to\R$ is \emph{short}, i.e.\ distance non-increasing.

\begin{figure}[H]
\centering
\begin{tikzcd}[column sep=huge, row sep=large]
(\R^n,\ \dd(u,v)=\norm{u-v}) \arrow[r,"\norm{\cdot}"] & (\R,\ \dd(s,t)=|s-t|)
\end{tikzcd}\\[3mm]
\begin{tikzpicture}[x=1cm,y=1cm]
  \draw[vec] (0,0) -- node[above left]{$u$} (1.3,1.1);
  \draw[vec] (1.3,1.1) -- node[above right]{$v-u$} (2.9,0.55);
  \draw[vec,dashed] (0,0) -- node[below]{$v$} (2.9,0.55);
  \node at (4.3,0.6) {$\longmapsto$};
  \draw[axis] (5.3,0.3) -- (9.4,0.3);
  \foreach \p/\l in {5.7/0, 7.1/{\norm{u}}, 8.6/{\norm{v}}}
    {\fill (\p,0.3) circle (1.4pt) node[below=2pt,font=\scriptsize]{$\l$};}
  \draw[decorate,decoration={brace,amplitude=4pt}] (7.1,0.55) -- (8.6,0.55)
    node[midway,above,font=\scriptsize]{$\le\norm{u-v}$};
\end{tikzpicture}
\caption{The bridge. Applying the arrow $\norm{\cdot}$ collapses the vector triangle of
\emph{Line\'aris algebra} onto the absolute-value triangle of \emph{Kalkulus}:
$\bigl|\norm{u}-\norm{v}\bigr|\le\norm{u-v}$. Certified in
\lean{DiagramProofs.length\_short} and \lean{TriangleOverlap.lengthShort}.}
\label{fig:bridge}
\end{figure}
